\section{Polarization and intrinsic critical native path obstructions}
\label{sec:136-prime-native-interpolation}

\paragraph{Archive, inputs, and scope.}
The source category and its uniform bounds are cote136, physical
pp.~117 and119, footers116 and118, scans117 and119; the
common-source and homotopy questions are physical p.~121,
footer120, scan121.  We use the checked internal statements
ARCH-136-RELATIVE-CELL-PATHS,
ARCH-136-DERIVED-TAIL-COMPARISON,
ARCH-136-RELATIVE-COHERENT-MOORE, and
ARCH-136-TAIL-EXT-VANISHING.  The strict source and its rational
comparison use ARCH-136-COHERENT-PD-REALIZATION; the final low-degree
repair uses ARCH-136-COHERENT-LOW-REPAIR and AQ-136-RELATIVE-H1.
The displayed polarization identity
is proved here.  No additional outside theorem or preprint is an
input.  The normalization, free PD universal property, and finite
relative-prism cohomology are exactly those already used and checked
in the cited modules.

The first result makes the critical obstruction of
ARCH-136-RELATIVE-CELL-PATHS independent of the native interpolation
when the native integer is prime.  The second constructs prescribed
endpoint maps with a nonzero obstruction.  Their common source and
its specified low-degree repair are provably inadmissible.  No counterexample to admissible-domain
simplicial homotopy, to the all-arrow groupoid assertion, or to
integral connectedness is deduced.

\subsection{An integral polarization identity}

Let \(q\geq1\) be odd, let \(p\) be prime, and put
\(D=p(q+1)\).  In any special strictly commutative dg-PD algebra,
for an odd element \(a\) of form degree \(q\), set
\[
 \omega_p(a)=\bigl(a\gamma_{p-1}(da),\gamma_p(da)\bigr)
                           \in P_p(C)^{D-1},
 \qquad d_{P_p}(x,y)=(dx-py,-dy).
                                                               \tag{136.PI1}
\]
It is closed because
\(d(a\gamma_{p-1}(da))=p\gamma_p(da)\).

Take \(a,c\) of the same odd degree and write \(b=da\),
\(e=dc\).  For \(1\leq i\leq p-1\), put \(j=p-i\) and
\[
 \begin{aligned}
 X_i&=a c\,\gamma_{i-1}(b)\gamma_{j-1}(e),\\
 Y_i&=a\,\gamma_{i-1}(b)\gamma_j(e),\\
 Z_i&=c\,\gamma_i(b)\gamma_{j-1}(e),\\
 W_i&=\gamma_i(b)\gamma_j(e).
 \end{aligned}                                                \tag{136.PI2}
\]
Their form degrees are \(D-2,D-1,D-1,D\), respectively.
The special differential and divided-power product formulas give
\[
                dX_i=iZ_i-jY_i,\qquad
                dY_i=iW_i,\qquad dZ_i=jW_i.
                                                               \tag{136.PI3}
\]
Choose an integer \(A_i\) with \(A_i i\equiv1\pmod p\),
and define the integers
\[
       B_i=\frac{-A_i j-1}{p},\qquad
       C_i=\frac{A_i i-1}{p}.
                                                               \tag{136.PI4}
\]
They are integral because \(j=p-i\).  Then
\[
 d_{P_p}\bigl(A_iX_i,B_iY_i+C_iZ_i\bigr)
                                         =(Y_i+Z_i,W_i).
                                                               \tag{136.PI5}
\]
Indeed the first component has coefficients
\(-A_i j-pB_i=1\) and \(A_i i-pC_i=1\), and
\(B_i i+C_i j=-1\) gives the second component with the
required minus sign.  Expanding divided powers of \(b+e\)
therefore proves the integral identity
\[
 \omega_p(a+c)-\omega_p(a)-\omega_p(c)
   =d_{P_p}\sum_{i=1}^{p-1}
                 \bigl(A_iX_i,B_iY_i+C_iZ_i\bigr).
                                                               \tag{136.PI6}
\]
For \(p=2\), it is simply
\[
       \omega_2(a+c)-\omega_2(a)-\omega_2(c)
                             =d_{P_2}(ac,-a\,dc).
                                                               \tag{136.PI7}
\]
If \(c\) restricts to zero on a simplicial subspace, every
term of the primitive in (136.PI6) does so too.  If \(a,c\)
have weight \(m\), all its terms have weight \(pm\).
The argument needs primality only to invert \(i\) modulo
\(p\); it makes no assertion of failure for composite integers.

\paragraph{ARCH-136-PRIME-INTERPOLATION-INVARIANCE.}
Retain every hypothesis and the actual pushout of
ARCH-136-RELATIVE-CELL-PATHS, with \(p\) prime.  Fix its
endpoint maps and complete old boundary data.  Its obstruction
\(\kappa(\widetilde a)\), defined in (136.RC11), is independent
of the chosen native interpolation \(\widetilde a\).
In the critical dimension \(n=D-2\), it is consequently one
well-defined element of \(S^v/pS^v\), with the chosen prism
orientation.  It vanishes if and only if the prescribed endpoints
admit a one-interval strict path under the old algebra.  No extra
weight bound is needed beyond the original input-row guard and
\(M+v\geq W_C(D)\).

\begin{proof}
Write \((Z,Y)\) for the relative prism in (136.RC8).
Two native interpolations \(\widetilde a_0,\widetilde a_1\)
have the same restriction to \(Y\).  Their difference
\(c=\widetilde a_1-\widetilde a_0\) is a section of
\(C^{m,q}(Z,Y)\).  The prescribed parameter cochain \(\beta\)
on \(Y\) is the same for both interpolations, so use one and
the same coordinatewise lift \(\widetilde\beta\).
The difference of their obstructions is the relative class of
\(\omega_p(\widetilde a_1)-\omega_p(\widetilde a_0)\),
with coefficients multiplied by the augmentation
\(\epsilon:P_0[v]\to I_M[v]\).

Apply (136.PI6).  Its mixed primitive is relative to \(Y\).
After multiplying its zeroth resolution component by \(\epsilon\),
it is also a primitive in the full resolution total complex:
all next resolution components are zero and
\(\epsilon\partial_1=0\).  Thus the obstruction difference
is exactly the class of the relative native pair \(\omega_p(c)\).

ARCH-136-RELATIVE-COHERENT-MOORE applies to this section \(c\)
on the arbitrary pair \((Z,Y)\), with the same input row,
resolution, and shifted parameter weights.  It supplies a full
relative coherent primitive for \(\omega_p(c)\), including
both bottom relations.  Its cone class is zero.  This proves
independence.  The necessary-and-sufficient extension statement
in ARCH-136-RELATIVE-CELL-PATHS then gives the stated criterion
for existence of a path; independence means that testing any one
native interpolation is enough.  This argument makes no claim
that coherent endpoint maps themselves cannot be changed.
\end{proof}

\subsection{A genuine nonzero prescribed critical component}

\paragraph{ARCH-136-CRITICAL-NATIVE-PATH-COMPONENT.}
Fix an integer \(p\geq2\), an odd \(q\geq1\), and set
\(D=p(q+1)\), \(n=D-2\).  Over \(\mathbb Z\), there is
an actual native coherent source \(E\), with \(E_0=E_1=S\),
and two strict maps \(f_0,f_1:E\rightrightarrows B\) to the
standard admissible theory, such that both native inputs are
zero and both maps agree on the complete old boundary, but they
are not joined by a strict cotensor path under that boundary.
This remains true for any finite oriented interval path in place
of one edge.  The source has the pointwise rational coefficient
comparison, but fails horizontal admissibility in form degree zero
in every weight.
The nonexistence of these prescribed paths persists after an RH1
iteration retaining both maps, simply by precomposition.

\begin{proof}
Choose a positive weight \(m\) in which the standard form-\(q\)
row is horizontally acyclic; such weights exist by the inherited
admissibility of the standard theory.  Put \(v=pm\).
Take \(M\) with \(M+v\geq W_B(D)\), and a free resolution
\(P_\bullet\to I_M\) with the previously specified generator
weights.  Let \(X=\Delta[n]\), \(\partial X=\partial\Delta[n]\).
Augment the complete boundary presentation to \(S\), and form
\[
                  E=S\amalg_{\mathcal N(\partial X)}
                                             \mathcal N(X).
                                                               \tag{136.PI8}
\]
Thus every native boundary generator and every boundary parameter
is zero.  The source is the actual independent coherent block
on the relative sphere; \(n\geq2\), so its first two simplicial
terms are \(S\).  The rational triangular-coordinate proof
already checked for the coherent presentation gives its pointwise
rational coefficient comparison.

Let
\[
 \mathcal H_B=\operatorname{Tot}
       \operatorname{Hom}^{0}_{S,\mathrm{gr}}
                    (P_\bullet[v],B(X,\partial X)).
\]
The derived-tail comparison, the coefficient-tail Ext theorem,
and the cone exact sequence give
\[
 H^{D-2}(P_p(\mathcal H_B))
       \cong H^n(X,\partial X;S^v/pS^v)
       \cong S^v/pS^v.
                                                               \tag{136.PI9}
\]
The comparison is valid through the required adjacent degrees.
Choose a closed cone cochain \(\beta\) representing the
nonzero class of \(T^{[v]}\) modulo \(p\).  Its components
\((-V,-U)\) are a complete coherent parameter system for zero
native input.  Closedness includes the two literal bottom
relations.  They vanish on \(\partial X\), so the universal
property gives a strict map \(f_1:E\to B\).  Let \(f_0\)
be the augmentation, followed by coefficients in \(B\).
Both maps send \(a,b\) to zero and preserve all of \(S\).

Consider first a single relative prism
\[
 Z=X\times\Delta[1],\qquad
 Y=(\partial X\times\Delta[1])\cup(X\times\partial\Delta[1]).
\]
The prescribed endpoint/boundary parameter cochain is zero on
\(\partial X\times\Delta[1]\) and endpoint zero, and is
\(\beta\) at endpoint one.  For the zero native interpolation,
its obstruction is, up to the orientation sign, the connecting
image of \([\beta]\) from (136.PI9).  That connecting map is
an isomorphism onto
\[
                       H^{n+1}(Z,Y;S^v/pS^v),
                                                               \tag{136.PI10}
\]
so the obstruction is nonzero.  One explicit way to verify the
isomorphism is to insert
\(Y_0=(\partial X\times\Delta[1])\cup(X\times\{0\})\).
The pair \((Z,Y_0)\) contracts to zero, while
\((Y,Y_0)\) is \((X,\partial X)\) at endpoint one.
The long exact sequence of this triple gives precisely that
suspension isomorphism.  Naturality of the derived and cone
comparisons identifies it with the parameter connecting map.

Any other putative native interpolation is itself zero on all
of \(Y\), because the native endpoints and old boundary are
zero.  ARCH-136-RELATIVE-COHERENT-MOORE therefore makes its own
relative native cone class zero with these same bounds.  Hence
it cannot change the nonzero connecting obstruction.  This last
argument works for every integer \(p\geq2\); primality is not
needed in this special zero-endpoint example.

If the interval is a finite oriented edge path, repeat the
argument with that path in place of \(\Delta[1]\).  Row
extension, derived relative comparison, and coherent Moore
primitives apply to the resulting simplicial pairs.  The path
contracts to its initial vertex, so the same triple still has
the suspension isomorphism (136.PI10).  The native input of a
putative path is still relative to the whole prescribed boundary,
and its coherent cone class is zero by the same theorem.  Thus
no such path can exist.

Finally, the supplied RH1 theorem is applicable with the two
retained admissible targets \(B,B\), since \(E_0=S\), weight
zero is \(\mathbb Z[0]\), and the rational comparison has
just been verified.  It produces strict extensions of both maps
and the asserted low-degree repair.  A path between those extended
maps would restrict along \(E\to E^+\) to a path already
excluded.  No horizontal or higher integral admissibility of
\(E^+\) is inferred; in fact its horizontal admissibility fails.

To see the exact failure for both domains, all faces between
\(E_1=S\) and \(E_0=S\) are the identity.  Hence the normalized
horizontal complex in form degree zero has \(N_1=0\),
\(N_0=S^w\), and \(H_0=S^w\ne0\) in every weight \(w\).
The guarded RH1 construction has no repair to make in simplicial
degrees zero and one: there are no one-form cycles there and
the zero-form coefficient comparison is already exact.  Its
normalized repairs therefore begin in simplicial degree two,
so \((E^+)_0=(E^+)_1=S\) as well.  The same row-zero
calculation proves that \(E^+\) is inadmissible.
\end{proof}

\paragraph{Verification status.}
The author completed a whole saved-file reading, including the
integral polarization coefficients, relative full-resolution signs,
both literal bottom equations, the suspension connecting map, and
the exact nonadmissible-source scope.  Root and reviewer123 each
completed a whole saved-proof reading and passed both statements.
The explicit row-zero calculation explains the inadmissibility
of the two restricted domains.  No finite computation is used, and
neither restricted source is an admissible object.
